\documentclass{beamer}

\usetheme{metropolis}

\title{Fermat's Test: Fast, Elegant, and Slightly Untrustworthy}
\author{Siddharth \& Amit}
\date{\today}

\newcommand{\Z}{\mathbb{Z}}
\newcommand{\bld}[1]{\textbf{#1}}
\newcommand{\itl}[1]{\textit{#1}}
\newcommand{\N}{\mathbb{N}}
\newcommand{\Q}{\mathbb{Q}}
\newcommand{\R}{\mathbb{R}}
\newcommand{\eps}{\varepsilon}
\newcommand{\poly}{\mathrm{poly}}
\newcommand{\polylog}{\mathrm{polylog}}
\newcommand{\bigO}{\mathcal{O}}

\begin{document}

\maketitle
%------------------------------------------------
\section{Probabilistic Primality Testing}

\begin{frame}{From Heuristics to Probability}

\begin{itemize}
    \item Heuristic tests can fail without warning
    \item Probabilistic tests provide \bld{provable error bounds}
\end{itemize}

\vspace{1em}

\begin{block}{Key Idea}
We allow a small probability of error in exchange for efficiency
\end{block}

\end{frame}

%------------------------------------------------

\begin{frame}{General Structure}

\begin{enumerate}
    \item Randomly pick $a$
    \item Check a number-theoretic condition involving $a$ and $n$
    \item If it fails $\Rightarrow$ $n$ is composite
    \item Otherwise, repeat
\end{enumerate}

\vspace{1em}

\begin{block}{Terminology}
If $a$ proves compositeness, it is called a \itl{witness}
\end{block}

\end{frame}

%------------------------------------------------

\begin{frame}{Outcome}

\begin{itemize}
    \item If any test fails $\Rightarrow$ definitely composite
    \item If all tests pass $\Rightarrow$ \bld{probably prime}
\end{itemize}

\vspace{1em}

\begin{block}{Guarantee}
Error probability decreases exponentially with repetitions
\end{block}

\end{frame}
%------------------------------------------------
\section{Fermat Primality Test}

\begin{frame}{The Idea}

If $a^{n-1} \not\equiv 1 \pmod{n}$, then $n$ is composite.

\vspace{1em}

If $a^{n-1} \equiv 1 \pmod{n}$, then $n$ is \itl{probably prime}.

\vspace{1em}

\begin{block}{Definition}
If $n$ is composite but satisfies
\[
a^{n-1} \equiv 1 \pmod{n},
\]
then $n$ is a \bld{pseudoprime to base $a$}.
\end{block}

\end{frame}

%------------------------------------------------

\begin{frame}{Example}

\[
n = 341 = 11 \cdot 31
\]

\[
2^{340} \equiv 1 \pmod{341}
\]

\vspace{1em}

\begin{block}{Observation}
341 is composite but passes the test (base 2)
\end{block}

\end{frame}

%------------------------------------------------
\section{Fermat's Little Theorem}

\begin{frame}{Statement}

\begin{block}{Theorem}
Let $p$ be prime and $p \nmid a$. Then
\[
a^{p-1} \equiv 1 \pmod{p}.
\]
\end{block}

\end{frame}

%------------------------------------------------

\begin{frame}{Setup for Proof}

Consider:
\[
S = \{1,2,3,\dots,p-1\}
\]

Multiply each element by $a$:
\[
aS = \{a,2a,3a,\dots,(p-1)a\}
\]

\vspace{1em}

\begin{block}{Claim}
$aS\pmod{p}$ is a permutation of $S$
\end{block}

\end{frame}

%------------------------------------------------

\begin{frame}{Proof: Uniqueness}

Assume:
\[
ai \equiv aj \pmod{p}
\]

Then:
\[
a(i-j) \equiv 0 \pmod{p}
\]

Since $\gcd(a,p)=1$, $a$ is invertible:
\[
i \equiv j \pmod{p}
\]

Since $1 \le i,j \le p-1$, we get $i=j$.

\end{frame}

%------------------------------------------------

\begin{frame}{Proof: Non-Zero Elements}

For any $i \in S$:
\[
ai \not\equiv 0 \pmod{p}
\]

since neither $a$ nor $i$ is divisible by $p$.

\vspace{1em}

Thus, $aS$ consists of $p-1$ distinct nonzero residues.

\end{frame}

%------------------------------------------------

\begin{frame}{Proof: Final Step}

Multiply all elements:

\[
a \cdot 2a \cdots (p-1)a \equiv 1 \cdot 2 \cdots (p-1) \pmod{p}
\]

\[
a^{p-1}(p-1)! \equiv (p-1)! \pmod{p}
\]

Since $(p-1)! \not\equiv 0 \pmod{p}$:

\[
a^{p-1} \equiv 1 \pmod{p}
\]

\end{frame}

%------------------------------------------------
\section{Algorithm}

\begin{frame}{Fermat Test Algorithm}

\begin{enumerate}
    \item Pick random $a$ with $\gcd(a,n)=1$
    \item Compute $a^{n-1} \bmod n$
    \item If $\neq 1$ $\Rightarrow$ composite
    \item Else $\Rightarrow$ probably prime
\end{enumerate}

\end{frame}

%------------------------------------------------

\begin{frame}{Time Complexity}

\begin{itemize}
    \item Modular exponentiation: $\bigO(\log n)$ multiplications
    \item Each multiplication: $\bigO((\log n)^2)$
\end{itemize}

\[
\Rightarrow \bigO((\log n)^3)
\]

\vspace{1em}

With $k$ iterations:
\[
\bigO(k (\log n)^3)
\]

\end{frame}

%------------------------------------------------
\section{Limitations}

\begin{frame}{Pseudoprimes}

\begin{itemize}
    \item Some composites pass the test
    \item Rare, but they exist
    \item There are only 21853 pseudoprimes base 2 that are less than $2.5\times10^{10}$. This means that for $n$ upto $2.5\times10^{10}$, Fermat's test fails only $8.7\times{10^{-5}}\%$ of times, 
\end{itemize}

\end{frame}

%------------------------------------------------

\begin{frame}{Carmichael Numbers}

\begin{block}{Definition}
$n$ is Carmichael if:
\[
a^{n-1} \equiv 1 \pmod{n}
\quad \forall a \text{ with } \gcd(a,n)=1
\]
\end{block}

\vspace{1em}

\begin{block}{Implication}
Fermat test fails for \bld{all} bases
\end{block}

\end{frame}

%------------------------------------------------

\begin{frame}{Korselt's Criterion}

A composite $n$ is Carmichael iff:

\begin{itemize}
    \item $n$ is square-free
    \item For every prime $p \mid n$:
    \[
    p-1 \mid n-1
    \]
\end{itemize}

\end{frame}

%------------------------------------------------

\begin{frame}{Example: 561}

\[
561 = 3 \cdot 11 \cdot 17
\]

\[
\begin{aligned}
2 &\mid 560 \\
10 &\mid 560 \\
16 &\mid 560
\end{aligned}
\]

\vspace{1em}

\[
a^{560} \equiv 1 \pmod{561}
\]

for all $\gcd(a,561)=1$

\end{frame}

%------------------------------------------------
%------------------------------------------------

\begin{frame}{Carmichael Numbers are Not Rare}

\begin{block}{Theorem (Alford--Granville--Pomerance, 1994)}
There are infinitely many Carmichael numbers.
\end{block}

\vspace{1em}

\begin{itemize}
    \item Infinitely many composite numbers satisfy:
    \[
    a^{n-1} \equiv 1 \pmod{n} \quad \forall \gcd(a,n)=1
    \]
    \item Fermat test fails on \bld{every base} for such $n$
\end{itemize}

\vspace{1em}

\begin{block}{Implication}
Fermat's test is fundamentally unreliable in the worst case
\end{block}

\end{frame}

%------------------------------------------------

\begin{frame}{Generating Carmichael Numbers}

\begin{block}{Korselt's Criterion}
$n$ is Carmichael iff:
\begin{itemize}
    \item $n$ is square-free
    \item For every prime $p \mid n$:
    \[
    p - 1 \mid n - 1
    \]
\end{itemize}
\end{block}

\vspace{1em}

\textbf{Construction Idea:}
\begin{itemize}
    \item Choose distinct primes $p_1, \dots, p_k$
    \item Form $n = p_1 p_2 \cdots p_k$
    \item Ensure $p_i - 1 \mid n - 1$ for all $i$
\end{itemize}

\vspace{1em}

\begin{block}{Reality}
Generates Carmichael numbers, but requires search
\end{block}

\end{frame}

\begin{frame}{Conclusion}

\begin{itemize}
    \item Fast: polynomial time
    \item Simple: easy to implement
    \item Flawed: fooled by Carmichael numbers
\end{itemize}

\vspace{1em}

\begin{block}{Takeaway}
Useful as a filter, not a proof of primality
\end{block}

\end{frame}

\end{document}